\begin{abstract}
Personalized urban route planning requires reconciling heterogeneous user preferences---culinary exploration, scenic sightseeing, leisure pacing, or speed-run touring---with spatial, temporal, and budgetary constraints. Existing approaches either rely on classical optimization solvers that cannot interpret natural language intent, or single-LLM pipelines that hallucinate POIs and produce spatially incoherent routes. We present \textbf{CityFlow}, an expert-ensemble multi-agent framework that decomposes route planning across 11 domain-specific experts orchestrated by a LangGraph state machine. Inspired by the Mixture-of-Experts principle of conditional computation through specialized modules, CityFlow features: (1)~scene-aware expert dispatch that classifies user intent into five scenario types and dynamically activates only relevant experts; (2)~parallel wave-based expert execution across two dependency-ordered waves; (3)~a review-rework feedback loop evaluating proposals against six quality dimensions with algorithmic spatial verification; and (4)~algorithmic time smoothing that validates and corrects LLM-generated temporal allocations. Evaluated on 5 core scenarios covering Zhuhai (2000+ POIs), CityFlow achieves an average score of 7.4/10 with 5/5 pass rate. An ablation study demonstrates that expert-ensemble decomposition outperforms single-LLM generation by 1.2 points on average, with the largest improvement on food exploration scenarios (+2.0 points). We document a systematic optimization journey from a baseline of 65.2/100 to 81.2/100 across 100 test scenarios, identifying seven categories of failure modes and their resolutions, yielding three design principles for LLM-based planning systems: (1)~prompt engineering should precede architectural complexity; (2)~LLMs should propose, algorithms should validate; and (3)~evaluation methodology errors can masquerade as algorithm failures. The system delivers first results within 5 seconds via Server-Sent Events streaming.
\end{abstract}
